% ================================================================
% ==== FILE: content/m_spaces/m1.tex
% ================================================================

% ==============================
%  Ontology of Continua — Core
%  Meta-Space: M1
%  FULL MODULE — FINAL
% ==============================

\subsubsection{\texorpdfstring{$M_1$}{M_1} Übersicht}
\label{sec:m1-overview}

$M_1$ ist der Metaraum, der die Zulässigkeitsstruktur für die bereitstellt
eindimensionales Kontinuum $K_1$.
Im Gegensatz zu $M_0$, das nur Proto-Differenzen regelt, führt $M_1$ das ein
erste echte geometrische und dynamische Zulässigkeitsbedingungen:
\[
\Omega(K_1) \subseteq \Omega(M_1).
\]

$M_1$ ist der minimale Metaraum, in dem Folgendes wohldefiniert ist:

\begin{itemize}
    \item eine messbare Achse $A_1$;
    \item ein kontinuierliches Feld $\phi(x)$ über einem eindimensionalen Bereich;
    \item Protodynamik über Flüsse $J_1=(\partial_x \phi, \partial_t \phi)$;
    \item zulässige Energiefunktionen und -wirkungen;
    \item wohlgeformte Strukturspannung $T_1$ und eine sinnvolle Schwelle $\Theta_1$;
    \item die Existenz- und Kollapsregeln von $K_1$.
\end{itemize}

Somit ist $M_1$ die Metastruktur, die $K_1$ mathematisch und physikalisch möglich macht.

\subsubsection*{1. Struktur von \texorpdfstring{$\Omega(M_1)$}{\Omega(M_1)}}

\[
\Omega(M_1)
=
\Big\{
\phi : X \to V \ \big|\
\phi \in C^{0}(X,V),\
\partial_x \phi,\ \partial_t \phi\ \text{wohldefiniert in } H^1
\Big\}.
\]

$M_1$ erzwingt:

\begin{enumerate}
    \item \textbf{Kontinuität:}
    Felder müssen kontinuierliche Darstellungen zulassen. Diskontinuierliche Abbildungen liegen außerhalb von $\partial\Omega(M_1)$.

    \item \textbf{Sobolev-Regularität:}
    Erste Ableitungen müssen im schwachen Sinne existieren:
    \[
    \partial_x \phi,\ \partial_t \phi \in H^1.
    \]

    \item \textbf{Metrikkompatibilität:}
    Die zulässige metrische Struktur von $X$ und $V$ ist durch $M_1$ festgelegt und wird zur Berechnung von Energie und Spannung benötigt.

    \item \textbf{Topologische Zulässigkeit:}
    $X$ muss eine eindimensionale topologische Mannigfaltigkeit sein
    (offenes Intervall, Kreis, Halblinie usw.).
\end{enumerate}

Jeder Verstoß gegen diese Einschränkungen führt zu einer Konfiguration an der Grenze
$\partial\Omega(M_1)$, was die Existenz von $K_1$ unmöglich macht.

\subsubsection*{2. Zulässige Achsen in \texorpdfstring{$M_1$}{M_1}}

$M_1$ erlaubt genau eine Strukturachse:
\[
A_1 : X \to \mathbb{R},
\]
stellt die erste messbare Dimension dar.

Folgendes ist in $M_1$ verboten:

\begin{itemize}
    \item zusätzliche unabhängige Achsen;
    \item inkompatible nichtlineare Struktur für $A_1$;
    \item Schwingachsen ohne zulässige Dynamik.
\end{itemize}

Diese Verbote definieren den Umfang möglicher $K_1$-Entwicklungen.

\subsubsection*{3. Zulässige Potenziale und Energien}

$M_1$ erfordert die Existenz einer wohldefinierten Energiefunktion
im Einklang mit der universellen Operatorstruktur:
\[
E[\phi] = \int_X \left(
\frac{1}{2} |\partial_x \phi|^2 + V(\phi)
\right)\,dx,
\]
mit:

\begin{itemize}
    \item $V(\phi)$ nach unten begrenzt,
    \item Koerzitivfeldstärke zur Verhinderung eines Zusammenbruchs,
    \item Differenzierbarkeit zur Definition von Flüssen.
\end{itemize}

Zulässige Potentiale $P(M_1)$ sind genau diejenigen, die die oben genannten Bedingungen erfüllen.

\subsubsection*{4. Schwellen und strukturelle Spannungen}

$M_1$ definiert den Schwellenwert:
\[
\Theta_1 = \inf \{ T_1(\phi) : \phi \in \Omega(M_1) \},
\]
und das Strukturspannungsfunktional:
\[
T_1(\phi) =
\int_X \left(
|\partial_x \phi|^2 + W(\phi)
\right)\, dx,
\]
wobei $W$ Bedingungen zur Durchsetzung von Einschränkungen enthält.

Eine Verletzung von $T_1 < \Theta_1$ führt zum Zusammenbruch von $K_1$.

\subsubsection*{5. Zulässige Flüsse und Protodynamik}

Zulässige Flüsse sind:
\[
J_1 = (\partial_x \phi,\ \partial_t \phi),
\]
wobei $\partial_t\phi$ zum zulässigen Tangentenraum von $\Omega(M_1)$ gehört.

$M_1$ verbietet Flüsse, die:
\begin{itemize}
    \item Spannungserhöhung ohne zulässigen Ausgleich,
    \item verursachen nichtphysikalische Diskontinuitäten,
    \item verstößt gegen das Aktionsprinzip oder Erhaltungsbeschränkungen.
\end{itemize}

\subsubsection*{6. Zusammenbruch und Grenze von \texorpdfstring{$M_1$}{M_1}}

Eine Konfiguration liegt auf der Grenze $\partial\Omega(M_1)$, wenn:

\begin{itemize}
    \item-Kontinuität schlägt fehl,
    \item Sobolev-Regularität schlägt fehl,
    \item Energie wird unbegrenzt,
    \item $T_1 \ge \Theta_1$,
    \item $A_1$ wird nicht messbar.
\end{itemize}

Der Kollaps von $K_1$ tritt auf, wenn seine zulässigen Konfigurationen diese Grenze überschreiten.

\subsubsection*{7. Beziehung zu \texorpdfstring{$K_0$}{K_0} und \texorpdfstring{$K_2$}{K_2}}

$M_1$ ist der minimale Raum, der den Übergang $K_0 \to K_1$ über den Operator $\Psi_{0\to1}$ ermöglicht.

Ein Übergang $K_1 \to K_2$ ist nur zulässig, wenn:
\[
A_2 \in A(M_2),\qquad
\Omega(K_2) \subseteq \Omega(M_2),\qquad
T_1 \ge \Theta_{\text{dim}}.
\]

Somit regelt $M_1$:

\begin{itemize}
    \item Existenz und Stabilität von $K_1$,
    \item der zulässige Bereich für das erste dynamische Kontinuum,
    \item die strukturelle Möglichkeit höherdimensionaler Kontinua.
\end{itemize}
